\documentclass[11pt,a4paper]{article}
\usepackage[margin=15mm]{geometry}
\usepackage{times}
\usepackage{setspace}
\usepackage[hidelinks]{hyperref}
\usepackage{enumitem}
\setlist[itemize]{leftmargin=1.2em,itemsep=0.15em,topsep=0.25em,parsep=0pt}
\pagestyle{empty}
\setlength{\parindent}{0pt}
\setlength{\parskip}{0.35em}
\setstretch{1.05}

\begin{document}

{\Large\bfseries Comparing Three Numerical Ways to Move a Two-Link Robot Arm\par}
\vspace{0.15em}
{\small\textit{One-page project proposal}}

\paragraph{Base paper.}
D.~Cardona-Ortiz and G.~Arechavaleta,
``Trajectory optimization for highly articulated robots based on sparsity-free local direct collocation,''
\textit{International Journal of Applied Mathematics and Computer Science}, 2025.
\url{https://doi.org/10.61822/amcs-2025-0041}\\[0.2em]
In plain words: they plan motions of \emph{large} robots by turning physics into a numerical problem.
We reuse that idea on a \emph{small} two-joint arm on a laptop. We will not rebuild their full industrial code.

\paragraph{Initial idea.}
We want a simple on-screen arm (two rods, two joints) to go from pose~A to pose~B.
The arm exists only as equations. The computer can (1)~guess motor commands and step physics forward in time, like a video game;
(2)~split time into a few points and connect them with a simple rule; or
(3)~use the same split, but a more accurate connecting rule (the style used in the base paper).
We run all three on the \emph{same} start, goal, and time budget, then compare who hits the target, who obeys physics better, and who is slower.
A matplotlib animation will play the three arms side by side so the difference is visible.

\paragraph{What we will produce.}
\begin{itemize}
  \item Python code for the two-link arm and the three numerical recipes.
  \item A side-by-side animation (same clock, same target dot).
  \item A small table: miss distance, physics error, and runtime versus number of time points.
\end{itemize}

\paragraph{How this matches the course directions.}
\textit{Cross-domain application:} we take course tools (stepping through time, connecting points, solving equations) and use them on a robot-arm motion problem, not on a textbook-only example.\\[0.25em]
\textit{Exploratory and comparative study:} same arm and same task; three methods measured side by side on a simple model (not a claim of a new algorithm).

\paragraph{Out of scope.}
No real robot. No 6-joint industrial software from the base paper.

\paragraph{What the demo should look like (existing examples).}
Our animation will look like a two-rod drawing that moves, not a factory video.
Look-alikes (not our code):
\href{https://matplotlib.org/stable/gallery/animation/double_pendulum.html}{Matplotlib double-pendulum animation}
and
\href{https://github.com/rparak/2-Link_Manipulator}{a two-link arm in Python (rparak)}.
Our contribution is the three-method comparison, not a copy of those demos.

\end{document}
